\documentclass[../../../include/open-logic-section]{subfiles}

\begin{document}

\olfileid{sth}{z}{power}
\olsection{مجموعات القوى}

سنواصل بمسلّمة أخرى:

\begin{axiom}[مجموعات القوى]
لكل مجموعة~$A$، توجد المجموعة~$\Pow{A} = \Setabs{x}{x \subseteq A}$.
\[
	\forall A \exists P \forall x(x \in P \liff (\forall z \in x)z \in A)
\]
\end{axiom}

تبريرنا لذلك مباشر جدًا. افترض أن~$A$ تكونت في المرحلة~$S$. عندئذ
كانت جميع عناصر~$A$ متاحة قبل~$S$ (بموجب~\stagesacc). ولذلك، باستدلال
مماثل لتبريرنا للفصل، تكون كل مجموعة جزئية من~$A$ قد تكونت بحلول
المرحلة~$S$. ومن ثم تكون جميعها متاحة لكي تُكوّن في مجموعة واحدة في
أي مرحلة بعد~$S$. ونعلم أن مرحلة كهذه توجد، لأن~$S$ ليست المرحلة
الأخيرة (بموجب~\stagessucc). ولذلك توجد~$\Pow{A}$.

وفيما يأتي نتيجة حسنة لمسلّمة مجموعات القوى:

\begin{prop}\label{thm:Products}
لكل مجموعتين~$A,B$، يوجد حاصل ضربهما الديكارتي~$A \times B$.
\end{prop}

\begin{proof}
توجد المجموعة~$\Pow{\Pow{A \cup B}}$ بموجب مجموعات القوى
و\olref[sth][z][pairs]{prop:pairsconsequences}. ولذلك، بموجب الفصل،
توجد المجموعة:
\[
	C = \Setabs{z \in \Pow{\Pow{A \cup B}}}{(\exists x \in A)(\exists y \in B) z = \tuple{x, y}}.
\]
والآن، لكل~$x \in A$ و$y \in B$، توجد المجموعة~$\tuple{x,y}$ بموجب
\olref[sth][z][pairs]{prop:pairsconsequences}. وفوق ذلك، لما كان
$x,y \in A \cup B$، كان~$\{x\},\{x,y\} \in \Pow{A \cup B}$،
وكان~$\tuple{x,y} \in \Pow{\Pow{A \cup B}}$. ومن ثم~$A \times B = C$.
\end{proof}

في هذا البرهان تتفاعل مسلّمة مجموعات القوى مع الفصل، وهذا غير مفاجئ.
فلولا الفصل لما كانت مجموعات القوى مبدأ \emph{قويًا} جدًا. فالفصل
يخبرنا أي المجموعات الجزئية من مجموعة ما توجد، وبذلك يحدد مقدار «ضخامة»
كل مجموعة قوى.

\begin{prob}
بيّن، لكل مجموعتين~$A,B$: (i) أن مجموعة جميع العلاقات التي مجالها~$A$
ومداها~$B$ موجودة؛ و(ii) أن مجموعة جميع الدوال من~$A$ إلى~$B$ موجودة.
\end{prob}

\begin{prob}
لتكن~$A$ مجموعة، ولتكن~$\sim$ علاقة تكافؤ على~$A$. أثبت وجود مجموعة
فئات التكافؤ تحت~$\sim$ على~$A$، أي~$\equivclass{A}{\sim}$.
\end{prob}

\end{document}
